% Осень26!A12; нумерация внутри столбца без пустых строк.
% Источник: лекция 4, §14 (строка 454)
\ticket{10}{Теорема Ферма. Необходимое условие экстремума. Теорема Ролля}
\label{ticket:10}

\subsection*{Локальные экстремумы: определения, теорема Ферма, необходимое условие, теорема Ролля}

Новый параграф --- «экстремумы». Речь исключительно о \textbf{вещественнозначных} функциях \(f:D\subset\mathbb{R}^m\to\mathbb{R}\) (отображения в многомерные пространства отложены); \(D\) обычно открыто, точки --- внутренние.

\subsubsection*{Определения}

Пусть \(a\in\mathrm{Int}\,D\).

\begin{itemize}
\tightlist
\item
  \textbf{Локальный максимум:} \(\exists U(a)\subset D:\ \forall x\in U(a)\ f(x)\le f(a)\) (нестрогое --- поблизости могут быть точки с тем же значением).
\item
  \textbf{Строгий локальный максимум:} \(\forall x\in\dot U(a)\) (проколотая окрестность) \(f(x)<f(a)\). Вставка прокола в предыдущее определение смысла не меняет, а здесь --- даёт именно строгий знак.
\item
  \textbf{Локальный минимум} / \textbf{строгий локальный минимум}: аналогично с \(\ge\) (\(>\)).
\item
  \textbf{Локальный экстремум} --- максимум или минимум (строгий --- аналогично).
\end{itemize}

(Шутка лектора про «оценку получил, но не понял, положительная или отрицательная» --- про ситуацию «то ли максимум, то ли минимум, но мы ещё не разобрались».)

\subsubsection*{Теорема Ферма}

\begin{theorem}

\emph{Ферма, многомерная, минимальная версия.} Пусть \(f:D\subset\mathbb{R}^m\to\mathbb{R}\), \(x_0\in\mathrm{Int}\,D\) --- точка локального экстремума, и \(f\) дифференцируема в точке \(x_0\). Тогда для любого направления \(l\in\mathbb{R}^m\), \(|l|=1\): \[\frac{\partial f}{\partial l}(x_0)=0.\]

\end{theorem}

\textbf{Доказательство (в одну строчку --- сведение к одномерной теореме Ферма).} Возьмём прямую \(L=\{x_0+tl,\ t\in\mathbb{R}\}\) с направляющим вектором \(l\) и рассмотрим функцию одной переменной \[h(t):=f(x_0+tl).\] Так как \(x_0\) --- внутренняя точка, при малых \(t\) точки \(x_0+tl\in D\); \(h\) имеет локальный экстремум в точке \(t=0\). Кроме того, \(h\) --- композиция двух дифференцируемых отображений: \(f\) дифференцируема в \(x_0\), а сдвиг \(t\mapsto x_0+tl\) дифференцируем всюду; значит, \(h\) дифференцируема в нуле и \[h'(0)=\frac{\partial f}{\partial l}(x_0).\] По одномерной теореме Ферма («Ферма из \(\mathbb{R}^1\)», 1 семестр) \(h'(0)=0\).

\subsubsection*{Необходимое условие экстремума}

\textbf{Следствие 1.} В условиях теоремы Ферма \(F'(x_0)=O\) --- нулевой оператор, что равносильно \[\operatorname{grad}f(x_0)=\Bigl(\frac{\partial f}{\partial x_1}(x_0),\dots,\frac{\partial f}{\partial x_m}(x_0)\Bigr)=0,\] т.е. все частные производные равны нулю.

(Достаточно применить теорему к \(l=e_j\), \(j=1,\dots,m\); в теореме сказано сильнее --- производные по \textbf{всем} направлениям нулевые, но можно было бы ограничиться стандартными.)

\textbf{Практический смысл:} для поиска экстремумов дифференцируемой функции решаем систему \(m\) уравнений \[\frac{\partial f}{\partial x_1}=0,\quad\dots,\quad\frac{\partial f}{\partial x_m}=0\] с \(m\) неизвестными --- получаем «подозрительные» (стационарные) точки; «только в них экстремумы и могут содержаться».

\subsubsection*{Теорема Ролля (многомерная)}

\begin{theorem}

\emph{Ролль.} Пусть \(K\subset\mathbb{R}^m\) --- компакт (с непустой внутренностью), \(f\) непрерывна на \(K\), дифференцируема в \(\mathrm{Int}\,K\), и \(f\big|_{\partial K}=\mathrm{const}\). Тогда существует \(x_0\in\mathrm{Int}\,K\): \[\operatorname{grad}f(x_0)=0.\]

\end{theorem}

\begin{proof}

По теореме Вейерштрасса \(f\) достигает на компакте \(K\) наибольшего значения в точке \(x_0\) и наименьшего в точке \(x_1\).

\begin{itemize}
\tightlist
\item
  Если \textbf{обе} точки \(x_0,x_1\) лежат на границе \(\partial K\), то максимум = минимум = граничная константа, значит \(f\equiv\mathrm{const}\) на всём \(K\), и тогда \(\operatorname{grad}f=0\) в любой внутренней точке --- всё тривиально.
\item
  Иначе хотя бы одна из точек экстремума внутренняя; не умаляя общности \(x_0\in\mathrm{Int}\,K\), и по теореме Ферма \(\operatorname{grad}f(x_0)=0\).
\end{itemize}

\end{proof}
